% main3.tex -- H-main.  Section 3: modified candidates and bendiness; basic lemmas with proofs.
%======================================================================
\section{Modified candidates and the bendiness of a pair}\label{sec:virtual}
%======================================================================

A real candidate has no well-defined velocity once Harry reaches it, and the front of real candidates tears where
candidates are caught. We therefore follow \emph{modified} (``virtual'') candidates, which slow down inside the vision
disc instead of fleeing at full speed. They agree with the real candidates until the moment they are seen, so seeing
every modified candidate is as good as seeing every real one; and their front is a smooth curve for all time, which never
shrinks. Then we introduce the quantity that measures how bent the front is, and derive the few calculus facts about it
that the rest of the proof uses. Throughout, $P$ is a $1$-Lipschitz path with $P(0)=0$, $d=1$ and $0<\varepsilon<1$.

\subsection{The regularised field and virtual fronts}

\begin{definition}[Regularised field]\label{def:F}
Put $f(\rho):=1$ for $\rho\ge\varepsilon$ and $f(\rho):=(\rho/\varepsilon)(2-\rho/\varepsilon)$ for
$0\le\rho\le\varepsilon$. Define $F(x):=f(|x|)\,x/|x|$ for $x\ne0$, $F(0):=0$, $A:=DF$, and $a^*(\rho):=1/\rho$ for
$\rho\ge\varepsilon$, $a^*(\rho):=2/\varepsilon$ for $\rho<\varepsilon$.
\end{definition}

\begin{lemma}\label{lem:F}
\textup{(i)} $F\in C^1(\R^2)$, $|F|\le1$, and $F(x)=x/|x|$ for $|x|\ge\varepsilon$.
\textup{(ii)} With $\rho=|x|$, $\hat r=x/\rho$ and $\Pi=I-\hat r\hat r^{\mathsf T}$:
$A=f'(\rho)\hat r\hat r^{\mathsf T}+(f(\rho)/\rho)\Pi$ is symmetric, $0\le A\le a^*(\rho)I$, $A=\Pi/\rho$ for
$\rho\ge\varepsilon$, $A(0)=(2/\varepsilon)I$, and $f(\rho)/\rho-f'(\rho)=\rho/\max(\rho,\varepsilon)^2$.
\textup{(iii)} For every unit vector $e$: $|e^\perp\cdot A(x)e|\le|x\times e|/\max(\rho,\varepsilon)^2$.
\textup{(iv)} $(F(x)-F(y))\cdot(x-y)\ge0$ for all $x,y$, and $(x/|x|-y/|y|)\cdot(x-y)\ge0$ for $x,y\ne0$.
\end{lemma}

\begin{proof}
Write $F(x)=\sigma(|x|)x$ with $\sigma(\rho)=(2\varepsilon-\rho)/\varepsilon^2$ on $[0,\varepsilon]$ and $\sigma=1/\rho$ on
$[\varepsilon,\infty)$. Both pieces have value $1/\varepsilon$ and slope $-1/\varepsilon^2$ at $\varepsilon$, so $\sigma$ is
$C^1$ on $(0,\infty)$; near $0$, $F(x)=(2\varepsilon x-|x|x)/\varepsilon^2$ is $C^1$ with $DF(0)=(2/\varepsilon)I$. Also
$|F|=f=1-(1-\rho/\varepsilon)^2\le1$ inside the disc. This is (i). Differentiating, $DF=\sigma I+\rho\sigma'\hat r\hat
r^{\mathsf T}=(f/\rho)\Pi+f'\hat r\hat r^{\mathsf T}$, since $f=\rho\sigma$. For $\rho\ge\varepsilon$ the eigenvalues are
$1/\rho$ and $0$; for $\rho<\varepsilon$ they are $f/\rho=(2\varepsilon-\rho)/\varepsilon^2\in(1/\varepsilon,2/\varepsilon]$ and
$f'=2(\varepsilon-\rho)/\varepsilon^2\in[0,2/\varepsilon]$, and $f/\rho-f'=\rho/\varepsilon^2$. This is (ii). For (iii),
$e^\perp\cdot Ae=-(f/\rho-f')(e\cdot\hat r)(e^\perp\cdot\hat r)$ and $|e^\perp\cdot\hat r|=|x\times e|/\rho$; use (ii).
For (iv), $(F(x)-F(y))\cdot(x-y)=\int_0^1(x-y)\cdot A(y+s(x-y))(x-y)\,ds\ge0$; and
$(x/|x|-y/|y|)\cdot(x-y)=|x|+|y|-(x\cdot y)(1/|x|+1/|y|)\ge0$ because $x\cdot y\le|x||y|$.
\end{proof}

\begin{definition}[Virtual flow]\label{def:vflow}
Given a $C^1$ curve $x^\circ:[u_-,u_+]\to\R^2$ with $\partial_ux^\circ\ne0$ and a time $t_\circ\ge0$, the
\emph{virtual flow} of $P$ from $x^\circ$ at $t_\circ$ solves
\[
\partial_tx(u,t)=vF\bigl(x(u,t)-P(t)\bigr),\qquad x(u,t_\circ)=x^\circ(u),\qquad t\ge t_\circ .
\]
We write $x_u(t):=x(u,t)$, $r=r_u:=x-P$, $\rho:=|r|$, $\hat r:=r/\rho$, $w:=|\partial_ux|$,
$T:=\partial_ux/w=e^{i\psi}$ and $\phi:=\arg r-\psi$. \emph{Virtual candidates} $x_\theta(t)$ start at $t_\circ=0$ from
$x^\circ(\theta)=e^{i\theta}$. A \emph{working front} starts at $t_\circ=t_\td$ from a curve $x^0$ on $[0,U]$ with
$|\partial_ux^0|\equiv1$; a label $u$ of it is \emph{true with direction $\theta$} if $S_\theta$ exists on $[0,t_\td]$
and $x^0(u)=S_\theta(t_\td)$.
\end{definition}

\begin{lemma}[Basic properties of virtual flows]\label{lem:vflow}
For a virtual flow:
\textup{(a)} $x(\cdot,t)$ is $C^1$, $|\partial_tx|\le v$, $\partial_t\partial_ux=vA(r)\partial_ux$, and both $w$ and
the arclengths $l(a,b;t)=\int_a^bw\,du$ are non-decreasing in $t$ \textup(so $w\ge1$ on a working front\textup).
\textup{(b)} $w_t=vw\,T\cdot AT$, $T_t=v(AT-(T\cdot AT)T)$ and
$\psi_t=-v\rho\max(\rho,\varepsilon)^{-2}\sin\phi\cos\phi$. For $\rho\ge\varepsilon$: $w_t/w=v\sin^2\phi/\rho$ and
$\psi_t=-(v/\rho)\sin\phi\cos\phi$.
\textup{(c)} For fixed labels, $l$, $c$ and $E$ \textup(the latter two where $c>0$; Definition~\ref{def:twopt}\textup)
are $C^1$ in $t$, and their derivatives depend on $P$ only through $P(t)$.
\end{lemma}

\begin{proof}
The field $(t,x)\mapsto vF(x-P(t))$ is continuous, and $C^1$ in $x$ with $|D_x|\le2v/\varepsilon$ (Lemma~\ref{lem:F}).
By the standard theory \cite[Ch.~V]{Hartman}, solutions are global and unique, depend $C^1$ on the initial point, and
$\partial_ux$ solves the variational equation $\partial_t\partial_ux=vA(r)\partial_ux$. Hence
$\partial_t(w^2)=2v\,\partial_ux\cdot A\partial_ux\ge0$, which gives (a) and $w_t=vwT\cdot AT$. Then
$T_t=(\partial_ux)_t/w-\partial_ux\,w_t/w^2$ is the formula in (b), and $\psi_t=T^\perp\cdot T_t=vT^\perp\cdot AT$.
Lemma~\ref{lem:F}(iii) with $e=T$ (where $T\cdot\hat r=\cos\phi$, $T^\perp\cdot\hat r=\sin\phi$) gives $\psi_t$; for
$\rho\ge\varepsilon$, $T\cdot AT=T\cdot\Pi T/\rho=\sin^2\phi/\rho$. Part (c) follows since $w>0$ and $\partial_ux$ is
$C^1$ in $t$, and the time derivatives are integrals of $vA(x-P(t))\partial_ux$.
\end{proof}

\begin{lemma}[Virtual hits imply true sightings]\label{lem:seen}
If $\rho_\theta(t)\le\varepsilon$ for a virtual candidate, or $\rho_u(t)\le\varepsilon$ for a true label $u$ of a working
front with direction $\theta$ and some $t\ge t_\td$, then the true candidate $S_\theta$ is seen by time $t$.
\end{lemma}

\begin{proof}
Let $t'\le t$ be the first time at which $\rho\le\varepsilon$ (it exists by continuity; for a virtual candidate $t'>0$
since $\rho(0)=1>\varepsilon$). Up to $t'$ we have $\rho\ge\varepsilon$, hence $F(x-P)=(x-P)/\rho$: the virtual point
solves the true equation \eqref{eq:snitch}. By uniqueness for \eqref{eq:snitch} away from $H$, it coincides with
$S_\theta$ on $[0,t']$ (respectively on $[t_\td,t']$, where it starts from $S_\theta(t_\td)$). So $S_\theta$ exists up to
$t'$ and $|S_\theta(t')-P(t')|=\rho(t')\le\varepsilon$.
\end{proof}

\begin{lemma}[Chords never shrink]\label{lem:chords}
Let $X_a,X_b$ solve $X'=vV(X-P)$ on $[t_1,t_2]$, each with its own $V\in\{F,\,x\mapsto x/|x|\}$ \textup(for $x/|x|$,
as long as the point does not meet $P$\textup).
\textup{(a)} If both use $F$, then $|X_b-X_a|$ is non-decreasing; in particular, chords $c(a,b;t)$ of a virtual flow are
non-decreasing in $t$.
\textup{(b)} In general, if $|X_b(t_1)-X_a(t_1)|\ge2\varepsilon$, then $|X_b(t)-X_a(t)|\ge|X_b(t_1)-X_a(t_1)|$ on
$[t_1,t_2]$.
\end{lemma}

\begin{proof}
Put $p_a:=X_a-P$, $p_b:=X_b-P$ and $D:=X_b-X_a=p_b-p_a$. Then $(|D|^2)'=2v(V_b(p_b)-V_a(p_a))\cdot(p_b-p_a)$. If
$V_a=V_b$ this is $\ge0$ by Lemma~\ref{lem:F}(iv), which proves (a). For (b), note first: whenever $|D|\ge2\varepsilon$,
one of $|p_a|,|p_b|$ is at least $\varepsilon$, say $|p_b|$; there $V_b(p_b)=p_b/|p_b|=F(p_b)$, and $V_b(p_b)-V_a(p_a)$
is a difference of the single monotone field $F$ (if $V_a=F$) or $x/|x|$ (if $V_a=x/|x|$), so again $(|D|^2)'\ge0$.
Now if $|D(t_1)|>2\varepsilon$, let $\tau$ be the infimum of times with $|D|<|D(t_1)|$; if $\tau<t_2$, then
$|D|>2\varepsilon$ on a neighbourhood of $\tau$, where $|D|$ is non-decreasing, a contradiction. If
$|D(t_1)|=2\varepsilon$, replace $X_b(t_1)$ by $X_b(t_1)+\iota D(t_1)/|D(t_1)|$; the perturbed solutions exist on
$[t_1,t_2]$ for small $\iota>0$, keep distance $\ge2\varepsilon+\iota$ from $X_a$, and converge to $X_b$ as $\iota\to0$.
\end{proof}

\subsection{Bendiness: arc minus chord}

How do we measure ``nearly straight''? Take two points $x_a,x_b$ on the front. The arc between them has length $l$, the
straight segment has length $c$, and the difference $E=l-c\ge0$ vanishes exactly when the arc is straight. We use
$N:=\sqrt{2E}$, because of a simple identity (Lemma~\ref{lem:excess}(b)): $N^2=\int|T-e|^2\,ds$, where $e$ is the
direction of the chord. So $N/\sqrt l$ is the root-mean-square angle between the front and its chord. The bound the proof
maintains, $N\le M\lambda^{1/2+\gamma}$ with $M\lambda^\gamma=m(\lambda/R)^\gamma$ (here $\gamma=\frac23$ and
$m=\frac1{10}$ are locked numbers, $R$ is a fixed small length chosen in Section~\ref{sec:assembly}, and
$M:=mR^{-\gamma}$), says: \emph{on a piece of label
length $\lambda$, the front deviates from its chord by an average angle at most $m(\lambda/R)^{\gamma}$}, which tends
to zero as the piece shrinks. Figure~\ref{fig:bend} shows the quantities; whether a pair is ``near'' or ``far'' depends
on its distance from Harry measured in units of its own length.

\begin{figure}[!htbp]
\centering
\begin{tikzpicture}[scale=1.0,>=Stealth]
  % left: one pair
  \begin{scope}[shift={(0,0)},scale=1.15]
    \draw[primary!35,line width=1pt] (-0.5,-0.4) .. controls (0.2,0.5) and (1.3,1.3) .. (2.2,1.35)
       .. controls (3.1,1.4) and (3.9,0.7) .. (4.5,-0.3);
    \draw[primary,line width=2.2pt] (0.32,0.5) .. controls (0.9,1.1) and (1.5,1.33) .. (2.2,1.35)
       .. controls (2.9,1.37) and (3.3,1.0) .. (3.65,0.62);
    \fill[primary] (0.32,0.5) circle (2pt) node[left,font=\small] {$x_a$};
    \fill[primary] (3.65,0.62) circle (2pt) node[right,font=\small] {$x_b$};
    \draw[accent,thick,dashed] (0.32,0.5) -- (3.65,0.62) node[pos=0.62,below,font=\small] {chord $c$};
    \node[primary,font=\small] at (2.2,1.68) {arc $l$};
    \draw[->,thick,black!70] (0.75,0.96) -- ++(0.55,0.45) node[above,font=\footnotesize] {$T$};
    \draw[->,thick,black!70] (0.75,0.96) -- ++(0.75,0.027) node[right,font=\footnotesize] {$e$};
    \node[font=\small,align=center] at (2.0,-0.8) {bendiness: $E=l-c$, \ $N=\sqrt{2E}$,\\[2pt]
      $N^2=\int_a^b|T-e|^2\,ds$};
  \end{scope}
  % right: near and far pairs relative to Harry on a leg
  \begin{scope}[shift={(6.4,0)},scale=1.0]
    \draw[primary!35,line width=1pt] (0,0) .. controls (2.0,0.18) and (3.8,0.30) .. (6.0,0.55)
      .. controls (7.0,0.68) and (8.0,0.88) .. (9.0,1.12);
    \coordinate (Aj) at (0.5,0.04);
    \coordinate (P) at (1.1,0.09);
    \coordinate (Aj1) at (1.7,0.15);
    \draw[accent,very thick] (Aj) -- (P);
    \draw[accent,thick,dashed] (P) -- (Aj1);
    \fill[accent] (P) circle (2.4pt) node[above=2pt,font=\small] {Harry};
    \fill[primary] (Aj) circle (1.6pt);
    \fill[primary] (Aj1) circle (1.6pt);
    \node[font=\footnotesize,primary] at (0.5,-0.25) {$x_{a_j}$};
    \node[font=\footnotesize,primary] at (1.85,-0.15) {$x_{a_{j+1}}$};
    % near pair
    \draw[primary,line width=2.2pt] (2.3,0.2) .. controls (2.7,0.23) and (3.1,0.26) .. (3.4,0.29);
    \node[font=\footnotesize,primary,align=center] at (2.85,0.75) {near pair\\(gap $<K_f\lambda$)};
    % far pair
    \draw[primary,line width=2.2pt] (6.6,0.63) .. controls (6.9,0.67) and (7.2,0.72) .. (7.5,0.79);
    \node[font=\footnotesize,primary,align=center] at (6.9,1.4) {far pair\\(gap $\ge K_f\lambda$)};
    \draw[<->,muted] (0.5,-0.6) -- (6.6,-0.6) node[midway,below,font=\footnotesize] {gap $g=a-a_j$};
    \draw[<->,muted] (6.6,-0.6) -- (7.5,-0.6) node[midway,below,font=\footnotesize] {$\lambda$};
  \end{scope}
\end{tikzpicture}
\caption{Left: the bendiness of a pair of points is the arc length minus the chord length; it is zero exactly for a
straight piece. Right: Harry walks the front in short legs; a pair is \emph{near} if its distance ahead of Harry is less
than ten times its own length, and \emph{far} otherwise (legs, the base point $x_{a_j}$ and the gap $g$ are defined in
Section~\ref{sec:chase}).}
\label{fig:bend}
\end{figure}

\begin{definition}[Two-point quantities]\label{def:twopt}
For a regular $C^1$ curve $u\mapsto x_u$ (arclength element $ds=w\,du$) and labels $a<b$:
\[
l(a,b):=\int_a^bds,\;\; c(a,b):=|x_b-x_a|,\;\; E:=l-c\ge0,\;\; N:=(2E)^{1/2},\;\; e=e_{ab}:=\frac{x_b-x_a}{c}\;\;(c>0).
\]
Set $\varsigma:=T\cdot e^\perp$, $q:=e^\perp\cdot A(r)e$, and, for a unit vector $\hat u$ and an interval $J$,
$V_{\hat u}(J):=\int_J|T-\hat u|^2\,ds$ and $\mathbf v_J:=x_{b'}-x_{a'}=c(J)e_J$ for $J=(a',b')$.
\end{definition}

\begin{lemma}[Excess and curvature]\label{lem:excess}
For a regular $C^1$ curve and an interval $J=(a,b)$:
\textup{(a)} $E(a',b')\le E(a,b)$ for $a\le a'<b'\le b$.
\textup{(b)} $V_{\hat u}(J)=2[l(J)-(x_b-x_a)\cdot\hat u]=N(J)^2+c(J)|e_J-\hat u|^2$ \textup(if $c(J)>0$\textup). Thus
$\int_J|T-e_J|^2\,ds=N(J)^2$, $V_{\hat u}$ is additive over adjacent intervals, and
$|\hat u\times(x_b-x_a)|\le(l\,V_{\hat u})^{1/2}$.
\textup{(c)} If the tangent angle is $K$-Lipschitz in arclength on $J$ and $Kl\le\pi$, then $N\le(K/\sqrt{12})\,l^{3/2}$.
\end{lemma}

\begin{proof}
(a) $l(a,b)=l(a,a')+l(a',b')+l(b',b)$ while $c(a,b)\le l(a,a')+c(a',b')+l(b',b)$. (b) $|T-\hat u|^2=2-2T\cdot\hat u$ and
$\int_JT\,ds=x_b-x_a=c\,e_J$, so $V_{\hat u}=2l-2c\,e_J\cdot\hat u=2(l-c)+c(2-2e_J\cdot\hat u)$. Also
$\hat u\times(x_b-x_a)=\int_J\hat u\times(T-\hat u)\,ds$; apply Cauchy--Schwarz. (c) With $\psi_m$ the tangent angle at
the midpoint, $|\psi(s)-\psi_m|\le K|s-l/2|\le\pi/2$, so
$c\ge\int_0^l\cos(\psi-\psi_m)\,ds\ge\int_0^l\cos(K|s-l/2|)\,ds=\frac2K\sin\frac{Kl}2\ge l-\frac{K^2l^3}{24}$.
\end{proof}

The next lemma says that a short piece inside a nearly straight longer piece points in nearly the same direction.

\begin{lemma}[Split tilt]\label{lem:splittilt}
For intervals $J\subset I$, $J\ne I$, with $c(I),c(J)>0$, put $l_c:=l(I)-l(J)$ and $Z:=|\mathbf v_J\times e_I|$. Then
\textup{(i)} $Z^2\le N(I)^2\,l(J)\,l_c/l(I)$;
\textup{(ii)} if $N(I)^2<4c(J)$, then $|e_J-e_I|\le(Z/c(J))\bigl(1-N(I)^2/(4c(J))\bigr)^{-1/2}$.
\end{lemma}

\begin{proof}
(i) $\mathbf v_J\times e_I=\int_J(T-e_I)\times e_I\,ds$, and also $=-\int_{I\setminus J}(T-e_I)\times e_I\,ds$ since
$\mathbf v_I\times e_I=0$. Cauchy--Schwarz gives $Z^2\le l(J)\int_J|T-e_I|^2$ and $Z^2\le l_c\int_{I\setminus J}|T-e_I|^2$.
Divide by $l(J)$ and $l_c$, add, and use Lemma~\ref{lem:excess}(b): $Z^2\,l(I)/(l(J)l_c)\le N(I)^2$.
(ii) With $\mathbf d:=e_J-e_I$, $|\mathbf d|^2(1-|\mathbf d|^2/4)=|e_J\times e_I|^2=Z^2/c(J)^2$, and
Lemma~\ref{lem:excess}(b) gives $c(J)|\mathbf d|^2\le V_{e_I}(J)\le V_{e_I}(I)=N(I)^2$.
\end{proof}

\subsection{How the bendiness changes in time}

The next lemma computes the time derivative of the bendiness of a fixed pair. Its last part is the key estimate for
pairs far from Harry: there the field $x/|x|$ is smooth on the scale of the pair, and the bendiness can grow only
through the sideways offset $p_a$ of Harry from the chord.

\begin{lemma}[Evolution of $l$, $c$ and $E$]\label{lem:Eprime}
For a virtual flow, fixed labels $a<b$ and $c(a,b;t)>0$ \textup(integrals over the arc from $x_a$ to $x_b$\textup):
\begin{enumerate}[(i)]
\setlength{\itemsep}{1pt}\setlength{\parskip}{0pt}
\item $l'=v\int T\cdot AT\,ds$ and $c'=v\int e\cdot AT\,ds$.
\item $E'=v\int(T-e)\cdot A(T-e)\,ds-\frac v2\int(e\cdot Ae)|T-e|^2\,ds+v\int\varsigma q\,ds$, where
$\int|T-e|^2\,ds=N^2$, $\int\varsigma\,ds=0$ and $|\varsigma|\le|T-e|$.
\item For every constant $\bar q$: $E'\le v\int a^*(\rho)|T-e|^2\,ds+v\int\varsigma(q-\bar q)\,ds$; also
$E'\le vN\|AT\|_{L^2(ds)}$.
\item \textup(Far inequality\textup) Let $\underline\rho:=\min_{s\in[a,b]}\rho_s$ and $p_a:=r_a\times e$. If
$\underline\rho\ge\varepsilon$, then
\[
E'\le\frac v{\underline\rho}\Bigl(1+\frac l{\pi\underline\rho}+\frac{2l^2}{\pi^2\underline\rho^2}\Bigr)N^2
+\frac{2vl^{3/2}|p_a|N}{\pi\underline\rho^3},\qquad
l'\le\frac v{\underline\rho}\Bigl[N+\frac{l^{1/2}|p_a|+lN/\pi}{\underline\rho}\Bigr]^2 .
\]
\end{enumerate}
\end{lemma}

\begin{proof}
(i) $l'=\int w_t\,du$ with Lemma~\ref{lem:vflow}(b); $c'=e\cdot\partial_t(x_b-x_a)=e\cdot\int_a^bvA\,\partial_ux\,du$.
(ii) $E'=v\int(T-e)\cdot AT$ ($A$ is symmetric). Write $(T-e)\cdot AT=(T-e)\cdot A(T-e)+(T-e)\cdot Ae$ and expand
$T-e=-\frac12|T-e|^2e+\varsigma e^\perp$ (as $T\cdot e=1-\frac12|T-e|^2$). Also $\int\varsigma\,ds=e^\perp\cdot(x_b-x_a)=0$.
(iii) Drop the middle term of (ii) ($A\ge0$), use $A\le a^*I$, and subtract $\bar q\int\varsigma=0$. The second bound is
Cauchy--Schwarz applied to $E'=v\int(T-e)\cdot AT$.
(iv) On the arc $\rho\ge\varepsilon$, so $A=\Pi/\rho$; the first term of (ii) is at most $vN^2/\underline\rho$. For the
last term use arclength $\sigma\in[0,l]$ from $x_a$ and coordinates $y=y_1e+y_2e^\perp$; then $q=Q(r)$ with
$Q(y)=-y_1y_2/|y|^3$, and a direct computation gives $|\partial_{y_2}Q|\le|y|^{-2}$ and $|\partial_{y_1}Q|\le2|y_2|/|y|^3$.
Let $\delta(\sigma):=\int_0^\sigma\varsigma$; then $\delta(0)=\delta(l)=0$, $|y_2(\sigma)|=|r_a\cdot e^\perp+\delta|\le|p_a|+|\delta|$,
and by Wirtinger's inequality \cite[Thm.~257]{HLP} $\int\delta^2\le(l/\pi)^2\int\varsigma^2\le(l/\pi)^2N^2$, hence
$\int|\delta|\le l^{3/2}N/\pi$. Integrating by parts (at a fixed time $dr/d\sigma=T$),
\[
\int\varsigma q\,d\sigma=-\int\delta\,\bigl(\partial_{y_1}Q\,(T\cdot e)+\partial_{y_2}Q\,\varsigma\bigr)\,d\sigma
\le\int|\delta|\frac{2(|p_a|+|\delta|)}{\underline\rho^3}+\int\frac{|\delta||\varsigma|}{\underline\rho^2}
\le\frac{2|p_a|l^{3/2}N}{\pi\underline\rho^3}+\frac{2l^2N^2}{\pi^2\underline\rho^3}+\frac{lN^2}{\pi\underline\rho^2},
\]
which gives the bound on $E'$. For $l'$: $T\cdot AT=\sin^2\phi/\rho$ with $|\sin\phi|=|T\times r|/\rho$ and
$|T\times r|\le|T-e|\rho+|e\times r|=|T-e|\rho+|y_2|$; so $|\sin\phi|\le|T-e|+(|p_a|+|\delta|)/\underline\rho$, and
Minkowski's inequality gives $(\int\sin^2\phi)^{1/2}\le N+(l^{1/2}|p_a|+lN/\pi)/\underline\rho$.
\end{proof}

\begin{lemma}[Duhamel inequality]\label{lem:duhamel}
If $c(a,b;t_1)>0$ and $\underline\rho(t)\ge\varepsilon$ on $[t_1,t_2]$, then with the \emph{rate}
$\mathcal A:=(v/\underline\rho)(1+l/(\pi\underline\rho)+2l^2/(\pi^2\underline\rho^2))$ and the \emph{forcing}
$\mathcal B:=2vl^{3/2}|p_a|/(\pi\underline\rho^3)$,
\[
N(t_2)\le e^{\int_{t_1}^{t_2}\mathcal A}N(t_1)+\int_{t_1}^{t_2}e^{\int_t^{t_2}\mathcal A}\mathcal B(t)\,dt .
\]
If moreover $\underline\rho\ge Kl$ on $[t_1,t_2]$, then $\mathcal A\le\mathcal A^*(K)\,v/\underline\rho$, where
$\mathcal A^*(K):=1+1/(\pi K)+2/(\pi K)^2$.
\end{lemma}

\begin{proof}
By Lemma~\ref{lem:chords}(a), $c\ge c(t_1)>0$ on $[t_1,t_2]$, so $E$ is $C^1$ there (Lemma~\ref{lem:vflow}(c));
$\mathcal A$ and $\mathcal B$ are continuous ($\underline\rho$ is a minimum of a jointly continuous function over a
compact label set). For $\iota>0$ put $N_\iota:=(N^2+\iota^2)^{1/2}$. By Lemma~\ref{lem:Eprime}(iv),
$N_\iota N_\iota'=E'\le\mathcal AN^2+\mathcal BN\le\mathcal AN_\iota^2+\mathcal BN_\iota$, so
$N_\iota'\le\mathcal AN_\iota+\mathcal B$. Integrate this linear inequality and let $\iota\to0$. The last claim is
$l/\underline\rho\le1/K$.
\end{proof}

